\documentclass[12pt,a4paper,oneside]{book}
\usepackage[utf8]{inputenc}
\usepackage[ngerman]{babel}
\usepackage[T1]{fontenc}

% Hauptdatei des Projektes. Von hier aus werden die anderen Seiten eingebunden
\usepackage{amsmath}
\usepackage{amsfonts}
\usepackage{amssymb}
\usepackage{xspace}
\usepackage[normalem]{ulem}
\usepackage{graphicx}
\usepackage{color, colortbl}
\usepackage[table]{xcolor}
\usepackage{nicefrac}

%%%%%%%%%%%%%%%%%%%%%%%%%%%%%%%%%%%%%%%%%%%%%%%%%%%%%%%%%%%%%%%%%%%%%%%%%%%%%%%%%%%%%%%%%%
% In diesem Abschnitt gibst du deine persönlichen Daten sowie den Kontext deiner Arbeit an
% das \xspace am Ende der Einträge muss vorhanden bleiben. Es regelt das Whitespace-
% handling
%%%%%%%%%%%%%%%%%%%%%%%%%%%%%%%%%%%%%%%%%%%%%%%%%%%%%%%%%%%%%%%%%%%%%%%%%%%%%%%%%%%%%%%%%%
% persönliche Daten
\newcommand{\vorname}{Vorname\xspace}
\newcommand{\nachname}{Nachname\xspace}
\newcommand{\emailadresse}{name@fb.hs-fulda.de\xspace}
\newcommand{\matrikelnummer}{Matr.Nr.\xspace}


% Arbeitsspezifische Daten
\newcommand{\titel}{Ein Titel\xspace}
\newcommand{\untertitel}{Ein  Untertitel\xspace}
\newcommand{\modulname}{Bachelor-Thesis\xspace}
\newcommand{\pruefer}{Ein Prof.\xspace}
\newcommand{\semester}{SoSe 2020\xspace}
\newcommand{\ortderarbeit}{Ein Ort\xspace}
%%%%%%%%%%%%%%%%%%%%%%%%%%%%%%%%%%%%%%%%%%%%%%%%%%%%%%%%%%%%%%%%%%%%%%%%%%%%%%%%%%%%%%%%%%

%Um einzelne Seite im Querformat zu verwenden
\usepackage{pdflscape}

% Damit Tabellen über die ganze Seitenbreite gehen
% https://texwelt.de/wissen/fragen/446/wie-bekomme-ich-eine-tabelle-in-textbreite
\usepackage{booktabs}
\usepackage{tabularx}
% Abstand zwischen Tabelle und Überschrift
% http://www.chemieonline.de/forum/showthread.php?t=101576
%\setlength{\belowcaptionskip}{0.3em}
\usepackage[center]{caption}
\newcommand*{\thead}[1]{\multicolumn{1}{c}{\bfseries #1}}
       
\usepackage{tikz}
\usepackage{pgfplots, pgfplotstable}

%Hurenkinder uns Schusterjungen
\clubpenalty10000
\widowpenalty10000
\displaywidowpenalty=10000

\usepackage{hyperxmp} % XMP-Daten fuer die PDF-Datei
%Gimmmick: Linked Kapitel und Inhaltsverzeichnis, sowie Referenzen
\usepackage[pdftex, pdfa]{hyperref}
\hypersetup{
    colorlinks,
    citecolor=black,
    filecolor=black,
    linkcolor=black,
    urlcolor=black,    
    pdftitle = {\titel},
    pdfauthor = {\vorname \nachname},
    pdfsubject = {\untertitel},
    pdfkeywords = {\titel \untertitel \matrikelnummer},
    pdflang = de,
    %bookmarks = true,
    pdfdisplaydoctitle = true,
    colorlinks = true,
    plainpages = false,
    %allcolors = black,
    hypertexnames = false,
    %pdfpagelabels = true,
    %hyperindex = true,
    unicode = true,
    pdfcaptionwriter = {\textsl{\vorname \nachname}},
    pdfcontactaddress = {Kattenstr. 9},
    pdfcontactcity = {Kassel},
    pdfcontactpostcode = {34119},
    pdfcontactcountry = {Deutschland},
    pdfcontactregion = {Hessen},
    pdfcontactemail = {\emailadresse},
    pdfcontactphone = {0176-21842232},
    pdfcontacturl = {http://www.hs-fulda.de},
    pdfmetalang = {de},
}

\usepackage{prettyref}
\usepackage{titleref}

%Schriftart

%\renewcommand*\rmdefault{cmdh}
%\usepackage{lmodern}
%\usepackage{mathptmx}
%\usepackage[scaled=.90]{helvet}
%\usepackage{courier}
%\usepackage{bera}
%\usepackage{boisik}

%Inhaltsverzeichnis
\usepackage[tocgraduated]{tocstyle}
\usetocstyle{allwithdot}
\setcounter{tocdepth}{3}

% Formatierungshilfen: http://wwws.htwk-leipzig.de/~myagovki/latex/formatierungshilfen/

\usepackage[ddmmyyyy]{datetime}
\renewcommand{\dateseparator}{.}


% Für Kopf und Fußzeile: https://esc-now.de/_/latex-individuelle-kopf--und-fusszeilen/?lang=en
\usepackage[
  headsepline, plainheadsepline,
  footsepline, plainfootsepline
]{scrlayer-scrpage}
\pagestyle{scrheadings}
\clearscrheadfoot

    
\ihead*{
   \itshape \leftmark
}
\ohead*{\vorname \nachname}
\ifoot*{\today}
\ofoot*{\pagemark}





% anderthalbfacher Zeilenabstand
\usepackage[onehalfspacing]{setspace}

%Geometry der Seite festlegen
\usepackage{geometry}
\geometry{
  left=2.5cm,
  right=2.5cm,
  top=2.5cm,
  bottom=2.5cm,
  headheight=33pt
}

%Für Gestaltung des Layouts
\usepackage{blindtext}

%Bilbliothek und Literatur einbinden
\usepackage[style=authoryear-ibid,natbib=true,backend=biber,sorting=nty]{biblatex}
\usepackage[babel,german=guillemets]{csquotes}
\defbibheading{head}{{\chapter*{Literaturverzeichnis}\addcontentsline{toc}{chapter}{Literaturverzeichnis}}}
\addbibresource{bib/literatur.bib}

%https://texwelt.de/fragen/19153/pdf-dateien-in-anhang-einbinden-aber-formatierung-des-vorhergehenden-dokuments-ubernehmen
\usepackage{pdfpages}

\newcommand{\zb}{z.\,B.\xspace}
\newcommand{\ua}{u.\,a.\xspace}
\newcommand{\oge}{o.\,g.\xspace}
\newcommand{\dah}{d.\,h.\xspace}
\newcommand{\bzw}{bzw.\xspace}
\newcommand{\bzgl}{bzgl.\xspace}
\newcommand{\ggf}{ggf.\xspace}
\newcommand{\bspw}{beispielsweise\xspace}

%https://stackoverflow.com/questions/3175105/writing-code-in-latex-document
\newcommand{\code}[1]{\texttt{#1}}

%https://tex.stackexchange.com/questions/115467/listings-highlight-java-annotations
\usepackage{listings}
\usepackage{inconsolata}

\usepackage{color}

\definecolor{pblue}{rgb}{0.13,0.13,1}
\definecolor{pgreen}{rgb}{0,0.5,0}
\definecolor{pgrey}{rgb}{0.46,0.45,0.48}
\definecolor{mauve}{rgb}{0.58,0,0.82}

\usepackage{listings}

\lstset{
  frame=t,
  language=Java,
  aboveskip=1mm,
  belowskip=1mm,
  columns=flexible,
  showspaces=false,
  showtabs=false,
  numbers=left,
  breaklines=true,
  showstringspaces=false,
  breakatwhitespace=true,
  commentstyle=\color{pgreen},
  keywordstyle=\color{pblue},
  stringstyle=\color{mauve},
  basicstyle={\small\ttfamily},
  breaklines=true,
  breakatwhitespace=true,
  tabsize=2,
  moredelim=[il][\textcolor{pgrey}]{$$}, %$$ 
  moredelim=[is][\textcolor{pgrey}]{\%\%}{\%\%}
}


% Silbentrennung
\hyphenation{Soft-ware-pro-jekt}
\hyphenation{Soft-ware}



%https://www.dante.de/events/dante2015/Programm/vortraege/vortrage-partosch2.pdf


\immediate\pdfobj stream attr{/N 3} file{res/AdobeRGB1998.icc}
\pdfcatalog{%
/OutputIntents [ <<
/Type /OutputIntent
/S/GTS_PDFA1
/DestOutputProfile \the\pdflastobj\space 0 R
/OutputConditionIdentifier (Adobe RGB 1998)
/Info(Adobe RGB 1998)
>> ]
}

\input{glyphtounicode.tex}
\input{glyphtounicode-cmr.tex}
\pdfgentounicode=1
\pdfobjcompresslevel=0
\pdfinclusioncopyfonts=1


\begin{document}

%!TEX root = ../hausarbeit.tex
\begin{titlepage}
\newgeometry{
  left=2.5cm,
  right=2.5cm,
  top=2.5cm,
  bottom=2cm
}
\begin{figure}
  \centering
  \includegraphics[width=1.0\textwidth]{res/hs_fulda_logo_ntr.jpg}
%   \caption{Meine Abbildung}
%   \label{fig:mein-bild}
\end{figure}
    \centering
    Hochschule Fulda \\    Fachbereich Elektrotechnik \\ Fachbereich Sozialwesen
    \vspace{1.5cm}

    {\Huge \bfseries \titel \par}
    {\Large \itshape \untertitel \par}
    \vspace{2.5cm}

    
    \modulname\\
    im Studiengang\\ Bachelor of Science: Sozialinformatik \\
    \vspace{0.7cm}

    \semester\\

    \vspace{3cm}
    Prüfer: \pruefer\\

    \vfill

% Bottom of the page
    vorgelegt von \\ \vorname \nachname \\ Matrikel-Nr.: \matrikelnummer \\ \emailadresse \\ %eingereicht am: \today
\restoregeometry
\end{titlepage}


\pagenumbering{roman}
\setcounter{page}{2}

\chapter*{Abstract}
\markboth{Abstract}{}
\label{sec:Abstract}

\blindtext


\setcounter{tocdepth}{2}
\tableofcontents

\newpage

\pagenumbering{arabic}

% Abstand nach Absätzen
\setlength{\parskip}{0.3cm}

% falls man lieber den kompletten Inhalt in einer Datei haben möchte, folgende Zeile auskommentieren und die Chapter-Zeilen (nächster Abschnitt) kommentieren
%\input{src/inhalt.tex}

% falls man lieber eine Datei pro Kapitel bevorzug, kann man dies auch einfach realisieren, indem man diese wie folgt einbindet
\chapter{Einleitung}
\markboth{Einleitung}{}
\label{sec:Einleitung}

\blindtext

\chapter{Grundlagen}
\markboth{Grundlagen}{}
\label{sec:Grundlagen}

\blindtext

\section{Unterkapitel}
\blindtext

\subsection{Hier ist im Inhaltsverzeichnis schluss}
\blindtext

\subsection{Weitere Subsection}
\blindtext

\section{Ein anderes}
\blindtext

\section{Letztes Unterkapitel}
\blindtext

\section{Das war gelogen}
\blindtext




\chapter{Fragestellung}
\markboth{Fragestellung}{}
\label{sec:Fragestellung}

\blindtext

\chapter{Vorgehensweise}
\markboth{Vorgehensweise}{}
\label{sec:Vorgehensweise}

\blindtext

\chapter{Fazit}
\markboth{Fazit}{}
\label{sec:Fazit}

\blindtext

% diese müssen dann natürlcih in dem Ordner vorliegen

\newpage

% nocite lässt alle Werke anzeigen, auch solche die nicht verlinkt sind.
%\nocite{*}
\printbibliography[heading=head]

\newpage

% Rudi will kein Abbildungsverzeichnis
\listoffigures
\addcontentsline{toc}{chapter}{\listfigurename}

\listoftables
\addcontentsline{toc}{chapter}{Tabellenverzeichnis}

\lstlistoflistings
\addcontentsline{toc}{chapter}{Listings}

\newpage

%!TEX root = ../hausarbeit.tex
% Erklaerung der Selbststaendigkeit.
% Footer wird in config.tex durch \AtEndDocument{\ofoot{}} angepasst
\chapter*{Erklärung der Selbständigkeit}
\markboth{}{}
\addcontentsline{toc}{chapter}{Erklärung der Selbständigkeit}

Ich versichere, dass ich die vorliegende schriftliche Prüfungsleistung selbständig verfasst, keine anderen als die angegebenen Quellen und Hilfsmittel verwendet habe und die Stellen, die anderen Werken im Wortlaut oder dem Sinn nach entnommen sind, im Text jeweils mit Quellenbelegen kenntlich gemacht habe. Die Arbeit ist noch nicht anderweitig für Prüfungszwecke vorgelegt worden.
\\
\\
\\
\\
\\

%\begin{minipage}[c]{15,4cm}
\parbox{6cm}{\centering  \ortderarbeit, \today\hrule\strut\centering\footnotesize Ort, Datum}
\hfill
\vspace{-2cm}
\parbox{6cm}{\centering {\color{white} X} \hrule\strut \centering\footnotesize Unterschrift}
%\end{minipage}



%\vspace{-2cm} %hochschieben
%\includegraphics[width=0.4\textwidth,heigth=12pt]{res/Unterschrift.png}
%\vspace{2cm} % wieder runter


\end{document}

